\section{Residual Class Under Revealed-Sacrifice Observation}
\label{sec:residual_class}

\subsection{Re-characterization}

A natural alternative approach would define the residual class as
capabilities failing benchmarkability (identity-individuated,
non-repeatable, non-communicable).  The revealed-sacrifice framing
produces a differently-shaped and operationally cleaner residual
class.

\begin{definition}[Residual Class Under Revealed Sacrifice]
\label{def:residual_class}
The residual class is the set of capabilities \emph{structurally
outside the reach} of the sacrifice channel---capabilities whose
exercise, by their nature, produces no sacrifice event even with
unlimited observation time:
\[
  \ResS = \{c \in P : \text{no sacrifice-observable path exists
  in principle}\}.
\]
$\ResS$ is a \emph{structural} property of the capability, not a
property of the observation window.  A capability is in $\ResS$
because the sacrifice channel \emph{cannot} reach it (purely
private exercise with no material footprint), not because the
channel has not yet observed it.  Examples: private contemplative
experience, purely internal cognitive operations with no external
inputs or time diversion.
\end{definition}

\begin{remark}[Structural versus evidential absence]
$\ResS$ is distinct from the set of capabilities with no sacrifice
evidence in a given window.  A capability outside $\ResS$
(tradable in principle) may have no sacrifice evidence in a
particular window simply because no trade happened to cover
it---this is the \emph{dormant} category in the gap decomposition
(Definition~\ref{def:gap_decomp}), not the residual class.
Dormant capabilities can be resolved by waiting for more trade
data; residual capabilities cannot.
\end{remark}

\begin{remark}[Conservatism of $\ResS$ classification]
Classifying a capability as residual is a structural claim that
should be conservative: if there is any plausible mechanism by
which exercise of~$c$ could produce a sacrifice event (even
indirectly---e.g., an artistic practice that requires purchasing
supplies), $c$ should be classified as outside $\ResS$.  The
residual class should be small by construction.
\end{remark}

\subsection{Relationship to the benchmarkability-based residual class}

\begin{proposition}[Residual Class Intersection]
\label{prop:residual_intersection}
Let $\ResB$ be the benchmarkability-based residual class
(capabilities failing benchmarkability under a direct-measurement
approach) and $\ResS$ be the revealed-sacrifice residual class
(Definition~\ref{def:residual_class}).  Then:
\begin{enumerate}
\item[(a)] $\ResB$ and $\ResS$ overlap but neither contains the
  other.
\item[(b)] Capabilities in $\ResB \setminus \ResS$:
  identity-individuated capabilities that nonetheless involve
  observable sacrifice.  Example: an artistic practice that is
  identity-individuated but involves purchasing supplies, paying
  for studio space, foregoing labour hours.  The time-sacrifice
  channel captures a lower bound even though the capability is
  identity-individuated.
\item[(c)] Capabilities in $\ResS \setminus \ResB$: capabilities
  that are benchmarkable (communicable, repeatable, externally
  verifiable) but are exercised purely internally with no trade
  event.  Example: a cognitive capability (mathematical reasoning,
  strategic planning) that the agent exercises entirely in private
  thought.
\item[(d)] Capabilities in $\ResB \cap \ResS$:
  identity-individuated capabilities exercised purely privately
  with no sacrifice event.  Example: private contemplative
  experience with no material input or time-sacrifice signal.
\end{enumerate}
\end{proposition}

\begin{proof}
Each part is established by exhibiting a capability with the
claimed membership, which proves non-containment in the relevant
direction.

\emph{(a) Neither class contains the other.}
Parts (b) and (c) below produce witnesses for
$\ResB \setminus \ResS \neq \emptyset$ and
$\ResS \setminus \ResB \neq \emptyset$ respectively, ruling out
either containment.

\emph{(b) $\ResB \setminus \ResS$.}
Consider the capability ``a specific artist's studio practice.''
Under the benchmarkability criterion of~\citep[%
Definition~2]{lasser2026poset}, this is identity-individuated: the
artefact sequence is non-repeatable and non-communicable in the
benchmark sense, so the capability lies in $\ResB$.  But the
practice consumes materials (paid-for supplies, rent on studio
space) and displaces outside-option labour hours; every session
produces a money- or time-sacrifice event
(Section~\ref{sec:two_channels}) that emits a valid
$\volRlower$ signal on the underlying bundle.  Thus the capability
is reachable in principle by the sacrifice channel and so
$\notin \ResS$.

\emph{(c) $\ResS \setminus \ResB$.}
Consider ``silent mental rehearsal of a chess opening.''  This is
a cognitive capability that is repeatable, communicable (the agent
could narrate the reasoning), and externally testable (by
presenting the resulting moves), so the capability is
benchmarkable and lies outside $\ResB$.  But exercised purely in
internal thought it produces no material input, no time-sacrifice
trace (rehearsal during idle cycles is not a wage-foregoing
event), and no counterparty; no sacrifice event is emitted in
principle.  Thus the capability lies in $\ResS$.

\emph{(d) $\ResB \cap \ResS$.}
Consider ``private contemplative experience with no material
footprint.''  This is identity-individuated (non-communicable
phenomenology, so $\in \ResB$) and emits no sacrifice event (no
money, no time beyond idle cycles, no counterparty, so
$\in \ResS$).  The intersection is witnessed and therefore
non-empty.
\end{proof}

\begin{remark}[The new residual class is the correct carve]
The prior residual class was defined by the \emph{measurement
apparatus} (what benchmarks can test).  The new residual class is
defined by the \emph{observation channel} (what trades reveal).
The observation-channel carve is the correct one for a paper about
what the framework can observe under its privacy commitments: the
question is not ``what could we test if we had access?'' but
``what do we actually see?''
\end{remark}

\subsection{Bounding the residual class}

\begin{proposition}[Active-Agent Floor Transfers]
\label{prop:oi_floor}
For an active agent~$k$ whose observational-individuation
capability $d_k$ lies within the observation perimeter $O$
(Appendix~\ref{app:exercise_indicator}), the observational
individuation of~\citep[Corollary~2.1]{lasser2026horizon}
provides a $\volR$ floor: if agent~$k$ is actively participating,
$d_k$ is continuously exercised, contributing at least
$\OIfloor(k)$ to $\volRexact(O)$.  This floor transfers
independently of the sacrifice channel---active participation
\emph{is} the exercise event.
\end{proposition}

\begin{remark}[Channel separation]
The OI floor contributes exclusively to $\volRexact(O)$ in the
combined diagnostic
(Appendix~\ref{app:exercise_indicator}):
\[
  \volRcomb = \volRexact(O) + \volRlower(P \setminus O).
\]
It does \emph{not} contribute to $\volRlower$ (the
sacrifice-based bound outside $O$).  Conflating the two channels
would double-count.
\end{remark}
